\chapter{Operacje na dowodach i warunki wyprowadzalności}\label{ch:derivability-proof}

\section{Dlaczego nie wystarczy sprawdzać pojedynczych dowodów?}
W rozdziale~\ref{ch:tarski-truth} wprowadziliśmy trzy warunki
Hilberta--Bernaysa--Löba, oznaczone (D1)--(D3).
Pierwszy z nich mówi, że do \emph{konkretnego} dowodu możemy
przypisać arytmetyczny dowód zdania o jego istnieniu. Pozostałe
dwa dotyczą natomiast \emph{dowolnych kodów dowodów}, także tych,
które przyjmują niestandardowe wartości w modelach teorii. Musimy więc
w samej arytmetyce uzasadnić poprawność operacji na kodach.

Przez cały rozdział ustalamy efektywnie aksjomatyzowalną teorię
$T\supseteq\mathrm{PA}$ w języku arytmetyki oraz \emph{konkretny}
predykat $\operatorname{Prf}_T(p,x)$ z Lematu~\ref{res:7:7}.
Kod $p$ zawiera dowód, świadectwa rozpoznania aksjomatów i
certyfikat poprawnego przebiegu jego weryfikatora. Zgodnie
z przyjętą tam konstrukcją $\operatorname{Prf}_T$ jest formułą $\Delta_0$.
Stąd $\operatorname{Bew}_T(x):=\exists p\,\operatorname{Prf}_T(p,x)$
jest formułą $\Sigma_1$.

\begin{uwaga}[Dlaczego własność $\Delta_0$ wymaga dodatkowego uzasadnienia?]
Z faktu, że pewna relacja jest pierwotnie rekurencyjna, \emph{nie wynika
bez dalszego argumentu}, że jej wykres można przedstawić formułą
$\Delta_0$ w samym języku $\{0,S,+,\cdot\}$. W niniejszym
predykacie kod $p$ zawiera nie tylko wiersze dowodu, lecz także
świadectwa weryfikacji, a ich składniki i ślady obliczenia są
ograniczone przez rozmiar kodu. Weryfikacja takiego
\emph{certyfikowanego} napisu sprowadza się do ograniczonych
kwantyfikatorów po jego częściach. Właśnie z tego dodatkowego
sposobu kodowania, a nie ze słowa „obliczalny”, pochodzi
przyjęta postać $\Delta_0$. Alternatywnie wystarcza reprezentacja
$\Sigma_1$ predykatu dowodu, jeżeli dla niej osobno udowodnimy
wymagane wewnętrzne operacje na kodach.
\end{uwaga}

Przy wyborze prezentacji $T$ ustalamy dodatkowo, że aksjomaty
$\mathrm{PA}$ otrzymują osobny, łatwo sprawdzalny rodzaj
certyfikatu. Jest to dozwolone: jeśli $T\supseteq\mathrm{PA}$
w sensie zbioru twierdzeń, dołączenie wszystkich aksjomatów
PA do wyliczenia aksjomatów $T$ nie zmienia dowodliwych zdań.
Dzięki oznaczeniu ich typu kod dowodu PA można przekształcić
w kod dowodu $T$ za pomocą funkcji pierwotnie rekurencyjnej.
Ta konwencja jest częścią \emph{wyboru standardowego
predykatu dowodu}; nie zakładamy, że każda formuła
rozpoznająca dowody spełnia analogiczne własności.

\begin{uwaga}[Liczba, numerał i kod składni]
Jeżeli $a=\ulcorner\varphi\urcorner$ jest \emph{liczbą naturalną},
to $\overline a$ jest \emph{termem} arytmetyki. Zapis
$\operatorname{Bew}_T(\overline a)$ jest zdaniem, a
$T\vdash\operatorname{Bew}_T(\overline a)$ stwierdzeniem
metamatematycznym, że istnieje jego dowód. Gdy w poniższych
wzorach występuje funkcja obliczeniowa $m(p,q)$, rozumiemy ją
jako skrót definicyjny dla reprezentującej ją formuły;
nie dopisujemy nowego symbolu funkcyjnego do
$\mathcal L_{\mathrm{ar}}$ bez definicji.
\end{uwaga}

\section{Kodowanie ciągów o zmiennej długości w PA}
Przy czytaniu kolejnych dowodów warto pamiętać zasadniczą różnicę:
możliwość zakodowania \emph{każdego jednego} skończonego ciągu
w metateorii nie daje automatycznie formuły, która
w $\mathrm{PA}$ zapewnia kod \emph{dla każdej długości}.
Potrzebny jest dowód indukcyjny.

\begin{lemat}[Dołączanie jednego wyrazu do kodu ciągu]
\label{lem:xi-append-beta}
Niech $\beta(a,b,i)$ oznacza funkcję odczytu z
Lematu~\ref{res:7:4}. Istnieje formuła arytmetyczna
opisująca dopisanie wyrazu $c$ do ciągu o długości $n$,
dla której $\mathrm{PA}$ dowodzi
\[
\forall a\,\forall b\,\forall n\,\forall c\;\exists a'\,\exists b'
\left(\left[\forall i<n\;\beta(a',b',i)=\beta(a,b,i)\right]
\land\beta(a',b',n)=c\right).
\]
Równość $\beta(u,v,i)=w$ jest skrótem formuły
$B(u,v,i,w)$ z Lematu~\ref{res:7:4}. W~PA istnienie i
jednoznaczność takiego $w$ wynikają z dzielenia z resztą;
nie traktujemy $\beta$ jako pierwotnego symbolu języka.
\end{lemat}
\begin{proof}
W $\mathrm{PA}$ mamy dzielenie z resztą, istnienie
wspólnych wielokrotności skończonej rodziny liczb i
chińskie twierdzenie o resztach; ich dowody uzyskuje się
z algorytmu Euklidesa i indukcji. Wyjaśnimy konkretną
konstrukcję, a nie tylko istnienie kodu.

Dla $i<n$ niech $r_i:=\beta(a,b,i)$, a $r_n:=c$.
Istnienie odczytów $r_i$ i ich jednoznaczność są dowodliwe
w PA, więc przez indukcję i rekurencję na skończonych listach
możemy utworzyć ich maksimum $R$ oraz liczbę
$b':=(n+1)!(R+1)$. Arytmetyka PA dowodzi, że $b'>r_i$
i że każda dodatnia liczba nie większa od $n+1$ dzieli
$(n+1)!$, a zatem również $b'$. Ustalamy
$m_i:=1+(i+1)b'$ dla $i\le n$. Liczby $m_i$ są parami
względnie pierwsze: wspólny dzielnik $m_i$ i $m_j$ dzieli
$(j-i)b'$ i jest względnie pierwszy z $b'$, a więc dzieli
$j-i$; ponieważ $j-i$ dzieli $b'$, musi być równy $1$.
Dodatkowo $r_i<m_i$, ponieważ $r_i\le R$ oraz
$m_i\ge 1+b'>R$.

Chińskie twierdzenie o resztach zapewnia $a'$ takie, że
$a'\equiv r_i\pmod {m_i}$ dla wszystkich $i\le n$.
Formalizacja nie ogranicza się do ustalonych numerałów:
prowadzimy w PA indukcję po liczbie $k$ już uwzględnionych
kongruencji. Jeśli $a_k$ spełnia pierwsze $k$ kongruencji,
a $P_k=\prod_{i<k}m_i$, to $P_k$ i $m_k$ są względnie
pierwsze (każdy $m_i$ jest względnie pierwszy z $m_k$).
W PA algorytm Euklidesa uzasadnia istnienie współczynników
Bézouta, a więc odwrotności $P_k$ modulo $m_k$.
Wybieramy resztę $0\le v<m_k$ spełniającą
$P_kv\equiv r_k-a_k\pmod{m_k}$ i kładziemy
$a_{k+1}:=a_k+P_kv$, redukując wynik modulo $P_km_k$.
Wtedy $a_{k+1}$ spełnia wszystkie pierwsze $k+1$
kongruencje. Istnienie odpowiednich reszt, współczynników
i iloczynów dla każdego $k\le n+1$ jest dowodliwe przez
indukcję PA, a nie przez nieskończoną listę instancji.
Otrzymujemy kod $a'$ dla \emph{dowolnego} $n$.
Warunki na reszty dają tezę z definicji funkcji $\beta$.
\end{proof}

\begin{uwaga}[Co było tu naprawdę istotne?]
Indukcja odbywała się w teorii $\mathrm{PA}$ po \emph{zmiennej}
liczbie kongruencji. W $Q$ wykazywaliśmy tylko konstrukcje
na ustalonych numerałach. Tego dodatkowego użycia indukcji
nie wolno przenosić do $Q$ bez dowodu.
\end{uwaga}

\begin{lemat}[Funkcje pierwotnie rekurencyjne w PA]
\label{lem:xi-pr-in-pa}
Dla każdej funkcji pierwotnie rekurencyjnej
$f:\N^k\to\N$ można wybrać formułę $F_f(\vec x,y)$
taką, że $\mathrm{PA}\vdash\forall\vec x\,\exists!y\,F_f(\vec x,y)$.
Arytmetyka Peana dowodzi także równań definiujących $f$.
\end{lemat}
\begin{proof}
Dowód prowadzimy indukcją po konstrukcji funkcji
pierwotnie rekurencyjnej. Funkcje bazowe (zero, następnik
i rzuty) są reprezentowane przez równania, których
istnienie i jednoznaczność wartości są bezpośrednie.
Jeżeli $f(\vec x)=g(h_1(\vec x),\ldots,h_r(\vec x))$,
to z dowodliwej totalności każdego $h_i$ otrzymujemy
wartości $z_i$, a z totalności $g$ wartość $y$;
jednoznaczność kolejnych wartości daje jednoznaczność $y$.

Pozostaje rekursja prosta:
$f(\vec x,0)=g(\vec x)$,
$f(\vec x,S(n))=h(\vec x,n,f(\vec x,n))$.
Nie używamy w formule symboli funkcyjnych $g,h$, lecz
ich reprezentacji $G(\vec x,z)$ i $H(\vec x,i,z,w)$.
Formuła $F_f(\vec x,n,y)$ stwierdza istnienie kodu ciągu
o długości $n+1$, którego wyraz początkowy $z_0$
spełnia $G(\vec x,z_0)$, ostatni jest równy $y$, a
$H(\vec x,i,z_i,z_{i+1})$ zachodzi dla każdego $i<n$.
Wszystkie odczyty $z_i$ wyrażamy formułą $B$; warunek
$\forall i<n$ jest ograniczony. Przy ustalonym $\vec x$
stosujemy schemat indukcji PA do zdania, że taki kod
istnieje dla długości $n+1$. Dla $n=0$ wystarcza
udowodniona totalność $G$. W kroku indukcyjnym totalność
$H$ dostarcza wyraz $z_{n+1}$, a
Lemat~\ref{lem:xi-append-beta} zapewnia jego dopisanie.
Dla dwóch kodów spełniających te same warunki indukcja
po $i\le n$ i jednoznaczność $G,H$ dają równość ich
$i$-tych wyrazów; w szczególności równe są wyniki $y$.
Tak PA dowodzi istnienia, jednoznaczności i równań
rekurencji dla formuły $F_f$.
\end{proof}

\section{Formalne łączenie dwóch dowodów}
Gdy mamy dowód $\varphi\to\psi$ oraz osobny dowód
$\varphi$, reguła eliminacji implikacji pozwala dopisać
wiersz $\psi$. Na kartce wygląda to prosto. Musimy
jednak upewnić się, że \emph{ta sama operacja działa
na kodach dowolnej długości}, i że $\mathrm{PA}$
może potwierdzić jej poprawność.

\begin{lemat}[Formalizacja modus ponens]\label{lem:xi-formal-mp}
Niech $T\supseteq\mathrm{PA}$ będzie efektywnie
aksjomatyzowalna. Dla ustalonych zdań $\varphi,\psi$
można skonstruować funkcję pierwotnie rekurencyjną
$m_{\varphi,\psi}:\N^2\to\N$, reprezentowaną w PA
przez formułę $M_{\varphi,\psi}(p,q,r)$, taką że
$\mathrm{PA}\vdash\forall p\,\forall q\,\exists!r\,
M_{\varphi,\psi}(p,q,r)$ oraz
\[
\mathrm{PA}\vdash\forall p\,\forall q\,\forall r\;\bigl(
 M_{\varphi,\psi}(p,q,r)\land
\operatorname{Prf}_T(p,\overline{\ulcorner\varphi\to\psi\urcorner})
\land\operatorname{Prf}_T(q,\overline{\ulcorner\varphi\urcorner})
\to\operatorname{Prf}_T(r,
\overline{\ulcorner\psi\urcorner})\bigr).
\]
Zapis $\operatorname{Prf}_T(m_{\varphi,\psi}(p,q),x)$
można stosować jedynie w rozszerzeniu definicyjnym PA;
w pierwotnym języku oznacza podstawienie wyniku opisanego
przez $M_{\varphi,\psi}$.
\end{lemat}
\begin{proof}
\emph{Konstrukcja algorytmu.} Dekodujemy z $p$ i $q$
kompletne dowody, łącznie z certyfikatami aksjomatów.
Następnie: (1) kopiujemy wszystkie wiersze pierwszego
dowodu; (2) kopiujemy wiersze drugiego, przesuwając
o stałą długość pierwszego dowodu numery wszystkich
odwołań do poprzednich wierszy; (3) dopisujemy
wiersz $\psi$, wskazując dwa końcowe wiersze jako
podstawy reguły eliminacji implikacji; (4) tworzymy
certyfikat akceptacji nowego kodu przez weryfikator.
Na niepoprawnych wejściach algorytm oddaje ustalony
kod zastępczy. Odczyty, przepisanie wierszy oraz
przesunięcia wskaźników wykonuje się przez rekursję
ograniczoną długością wejściowych list; są to operacje
pierwotnie rekurencyjne. Weryfikator $V_T$ z Lematu~\ref{res:6:5} jest
funkcją pierwotnie rekurencyjną. \emph{Wybieramy}
realizujący go program otrzymany konstrukcyjnie z jego
schematu pierwotnie rekurencyjnego, a nie dowolny
program obliczający tę samą funkcję. Dla funkcji
bazowych czas działania takiej realizacji ma oszacowanie
pierwotnie rekurencyjne; przy składaniu sumujemy
oszacowania wykonywanych podprogramów, a przy rekursji
prostej tworzymy ograniczenie kolejnych czasów przez
rekursję prostą po jej liczniku. Indukcja po \emph{tym
konkretnym schemacie} daje więc pierwotnie rekurencyjne
ograniczenie liczby kroków, a w konsekwencji pierwotnie
rekurencyjne wytwarzanie pełnego śladu. Tak powstaje
certyfikat akceptacji nowego dowodu. Dla arbitralnego
programu obliczającego funkcję pierwotnie rekurencyjną
czas działania nie musi być pierwotnie rekurencyjny;
wybór realizacji jest zatem istotny.

\emph{Dowód wewnątrz PA.} Nie wystarcza stwierdzenie,
że operacja działa na wszystkich zwykłych liczbach.
Definiujemy w arytmetyce następujący niezmiennik
dla $j$ kopiowanych wierszy: „każdy z pierwszych $j$
wierszy jest poprawnie uzasadniony, a wskaźniki
wcześniejszych przesłanek są przesunięte zgodnie
z długością poprzedniego fragmentu”. Przypadek $j=0$
jest pusty. W kroku indukcyjnym rozróżniamy, czy
kopiowany wiersz jest aksjomatem, czy wynikiem
reguły. W pierwszym przypadku zachowujemy
jego świadectwo pochodzenia z wyliczenia aksjomatów,
a w drugim wykorzystujemy założoną poprawność
jego przesłanek i przesunięcie ich wskaźników.
Ostatni nowy wiersz przechodzi test eliminacji implikacji.
Lematy o kodowaniu ciągów i funkcji pierwotnie
rekurencyjnych gwarantują, że PA potrafi zbudować
kod całej listy oraz ślad jej weryfikacji.
W ten sposób indukcja \emph{po zmiennej liczbie
wierszy} daje w PA dokładnie podaną implikację.
\end{proof}

\begin{przyklad}[Jak odczytać twierdzenie o kodach?]
Załóżmy, że $p$ koduje dowód „jeśli $A$, to $B$”,
a $q$ koduje dowód $A$. Funkcja $m(p,q)$ tworzy
\emph{nowy} kod zawierający oba fragmenty i dodatkowy
krok prowadzący do $B$. Lemat~\ref{lem:xi-formal-mp}
nie twierdzi, że $A$ jest prawdziwe wyłącznie dlatego,
że ktoś zapisał symbol $p$: mówi, że \emph{jeżeli oba
kody przechodzą formalne testy poprawności}, to
kod otrzymany przez ich połączenie również przechodzi
taki test. To odróżnia regułę od jej dowolnego użycia.
\end{przyklad}

\begin{twierdzenie}[Warunek (D2)]\label{thm:xi-d2}
Dla standardowego predykatu dowodliwości $\Box$
teorii $T\supseteq\mathrm{PA}$ zachodzi dla
wszystkich zdań $\varphi,\psi$:
\[
T\vdash\Box(\varphi\to\psi)\to
(\Box\varphi\to\Box\psi).
\]
\end{twierdzenie}
\begin{proof}
Niech $g=\ulcorner\varphi\urcorner$,
$h=\ulcorner\psi\urcorner$ oraz
$k=\ulcorner\varphi\to\psi\urcorner$.
Lemat~\ref{lem:xi-formal-mp} daje w $\mathrm{PA}$
implikację, że jeśli istnieją $p$ i $q$ spełniające
$\operatorname{Prf}_T(p,\overline k)$ i
$\operatorname{Prf}_T(q,\overline g)$, to
istnieje kod dowodu zdania o kodzie $h$.
Po zastosowaniu wprowadzenia kwantyfikatora egzystencjalnego
i twierdzenia o dedukcji otrzymujemy
\[
\mathrm{PA}\vdash
(\exists p\,\operatorname{Prf}_T(p,\overline k))\to
\bigl((\exists q\,\operatorname{Prf}_T(q,\overline g))\to
\exists r\,\operatorname{Prf}_T(r,\overline h)\bigr).
\]
Po rozwinięciu definicji $\Box$ jest to dokładnie (D2).
Ponieważ $T\supseteq\mathrm{PA}$, ten sam dowód
jest dostępny w $T$.
\end{proof}

\section{Jednolita wersja rozstrzygania formuł ograniczonych}
Z (D1) nie otrzymamy (D3) przez zwykłe podstawienie.
(D1) dotyczy rzeczywistego dowodu wybranego zdania,
a (D3) musi być \emph{jednym dowodem w PA} implikacji
zawierającej kwantyfikator egzystencjalny po kodach.

\begin{lemat}[Sformalizowana prawdziwość $\Delta_0$]
\label{lem:xi-delta-proofcompiler}
Niech $\delta(\vec x)$ będzie ustaloną formułą $\Delta_0$.
Istnieje funkcja pierwotnie rekurencyjna $b_\delta^+$,
która dla standardowego $\vec n$ z
$\mathfrak N\models\delta(\overline{\vec n})$
wytwarza kod dowodu zdania $\delta(\overline{\vec n})$.
Ponadto dla tego samego, \emph{jednego} algorytmu
konstrukcji $\mathrm{PA}$ dowodzi zależności jednolitej
\[
\forall\vec x\;\bigl(\delta(\vec x)\to
\operatorname{Prf}_T\bigl(b_\delta^+(\vec x),
\operatorname{inst}_\delta(\vec x)\bigr)\bigr),
\]
gdzie $\operatorname{inst}_\delta(\vec n)$ jest kodem
napisu $\delta(\overline{\vec n})$. Funkcja podstawiania
numerałów jest pierwotnie rekurencyjna na mocy
Lematu~\ref{res:6:5}; funkcja $b_\delta^+$ jest totalna,
także na wejściach, dla których $\delta$ jest fałszywa
(wtedy jej wynik nie musi być kodem poprawnego dowodu).
Funkcje w wyświetlonym wzorze są oznaczeniami
w rozszerzeniu definicyjnym PA. Dokładniej: jeśli
$F_\delta^+(\vec x,p)$ i $I_\delta(\vec x,z)$
reprezentują odpowiednio wykresy $b_\delta^+$ oraz
$\operatorname{inst}_\delta$, to w pierwotnym języku
arytmetyki twierdzenie przyjmuje postać
\[
\mathrm{PA}\vdash\forall\vec x\,\forall p\,\forall z\;\bigl(
[\delta(\vec x)\land F_\delta^+(\vec x,p)\land
 I_\delta(\vec x,z)]\to\operatorname{Prf}_T(p,z)\bigr).
\]
Istnienie i jednoznaczność wyników reprezentowanych funkcji
są dowodliwe w PA na mocy Lematu~\ref{lem:xi-pr-in-pa}.
\end{lemat}
\begin{proof}
W metateorii stosujemy indukcję \emph{po budowie
ustalonej formuły $\delta$}. Wzmocnimy niezmiennik
indukcji: konstruujemy równocześnie \emph{dwa} totalne
kompilatory pierwotnie rekurencyjne $b_\delta^+$
i $b_\delta^-$, dla których PA dowodzi
\[
\begin{aligned}
\delta(\vec x)&\to
 \operatorname{Prf}_T(b_\delta^+(\vec x),
              \operatorname{inst}_\delta(\vec x)),\\
\neg\delta(\vec x)&\to
 \operatorname{Prf}_T(b_\delta^-(\vec x),
              \operatorname{inst}_{\neg\delta}(\vec x)).
\end{aligned}
\]
Obie implikacje rozumiemy jako formuły rozszerzenia
definicyjnego PA; w pierwotnym języku zastępujemy każdy
występujący symbol funkcji przez reprezentującą go formułę
wykresu i używamy dowodliwej jednoznaczności jego wartości.
Wprowadzenie drugiego kompilatora jest konieczne: przy negacji
potrzebny jest dowód fałszywości podformuły, a nie tylko
dowód podformuły prawdziwej. Dla obu polaryzacji funkcje
pozostają totalne, także poza odpowiednią gałęzią warunku.
W każdym kroku uzasadniamy jednolicie poprawność programu
\emph{wewnątrz} PA.

\emph{Formuły atomowe.} Dla każdego podstawienia
numerałów do termów obliczamy ich wartości,
a następnie produkujemy skończone uzasadnienie
z aksjomatów dodawania, mnożenia i następnika.
Jeżeli atomem jest równość, rozstrzygamy ją
przez porównanie wynikowych numerałów; jeśli
jest relacją $<$, korzystamy z definicji porządku.
Dla nierównej pary wartości program tworzy dowód
nierówności, a dla nieprawdziwego porównania
program tworzy dowód jego negacji. W PA poprawność
tych \emph{obu} procedur dla wszystkich parametrów
wynika z indukcji po liczbie wykonywanych operacji,
a nie z osobnych sprawdzeń.

\emph{Spójniki.} Przy negacji zamieniamy role
kompilatorów dodatniego i ujemnego, dołączając
odpowiedni krok klasycznego rachunku zdań. Dla
koniunkcji prawdziwej łączymy dwa dowody składników;
przy jej fałszywości wystarczy dowód negacji jednego
składnika. Analogicznie dla alternatywy. Dla implikacji
$\eta\to\theta$ gałąź dodatnia korzysta z dowodu
$\neg\eta$ albo $\theta$, a gałąź ujemna łączy dowody
$\eta$ i $\neg\theta$ i stosuje odpowiednią tautologię
rachunku zdań. Wybór gałęzi opiera się na pierwotnie
rekurencyjnej procedurze rozstrzygania ustalonej formuły
$\Delta_0$. Wszystkie operacje na kodach dowodów są
pierwotnie rekurencyjne, a PA dowodzi ich poprawności
przez indukcję po długości przepisanych dowodów,
jak w Lemacie~\ref{lem:xi-formal-mp}.

\emph{Kwantyfikatory ograniczone.} Kwantyfikatory
z ograniczeniem $u\le t$ najpierw rozwijamy zgodnie
z Definicją~\ref{res:4:6} do przypadku $u<t$ oraz
jednej dodatkowej instancji $u=t$; dlatego wystarcza
rozpatrzyć ograniczenie ostre. Jeśli
$\delta(\vec x)=\exists u<t(\vec x)\,\eta(\vec x,u)$,
to dla danych $\vec n$ obliczamy
$k=t(\overline{\vec n})$ i badamy kolejno
$u=0,\ldots,k-1$. Jeśli istnieje świadek,
wybieramy najmniejszy i dołączamy odpowiedni
dowód instancji $\eta(\overline{\vec n},\overline u)$,
dowód nierówności $\overline u<t(\overline{\vec n})$
oraz wprowadzenie kwantyfikatora. Dla
$\forall u<t(\vec x)\,\eta(\vec x,u)$ w przypadku
prawdziwym łączymy $k$ dowodów poszczególnych
instancji. W przypadku fałszywym znajdujemy
konkretny indeks $j<k$ i dołączamy dowód
$\neg\eta(\overline{\vec n},\overline j)$.
Dla fałszywej formuły z kwantyfikatorem egzystencjalnym
łączymy natomiast dowody negacji \emph{wszystkich}
$k$ instancji. W ten sposób otrzymujemy także
kompilator ujemny dla obu typów kwantyfikatorów.
W każdym przypadku kompilator dołącza również dowód
równości $t(\overline{\vec n})=\overline k$,
skonstruowany z obliczenia wartości termu $t$.
Bez tego kroku dowód ograniczony przez numerał
$\overline k$ nie byłby jeszcze dowodem formuły
z oryginalnym ograniczeniem $t(\overline{\vec n})$.
Podstawienie równych termów zapewnia odpowiednie
przeniesienie ograniczenia.
Przejście od listy instancji do zdania z ograniczonym
kwantyfikatorem jest formalizowane w PA przez
indukcję po \emph{zmiennym} $k$. Dla konkretnego $k$
można wyprowadzić z aksjomatów PA zdanie
$u<\overline k\to(u=\overline0\lor\cdots\lor
u=\overline{k-1})$ (dla $k=0$ otrzymujemy sprzeczność
z $u<0$). Z dowodów $k$ instancji, po podstawieniu
równości w każdym przypadku, tworzymy dowód
$\forall u<\overline k\,\eta$; ten sam schemat
z dowodami negacji daje dowód
$\forall u<\overline k\,\neg\eta$, a przez prawa
logiki klasycznej --- zaprzeczenie ograniczonego
kwantyfikatora egzystencjalnego. Dla fałszywej formuły
ograniczonej uniwersalnie lub prawdziwej egzystencjalnie
wystarcza jeden świadek oraz odpowiednia reguła logiczna.
Operacja tworzenia dowodu rozkładu skończonego przedziału
ma kod wyznaczany pierwotnie rekurencyjnie z $k$.
Indukcja w PA formalizuje jej poprawność i sprawdzanie
całego wygenerowanego kodu dla zmiennego $k$;
nie zastępujemy jej nieskończoną rodziną dowodów
dla oddzielnych numerałów. Dla ustalonego $k$
wykorzystujemy także Lemat~\ref{res:4:4}, lecz sam ten lemat
nie zastępuje indukcji po zmiennym $k$ w PA.
Lemat~\ref{lem:xi-append-beta} gwarantuje możliwość
wydłużania kodu listy. W ten sposób dowód jednolity
nie jest tylko nieskończoną rodziną oddzielnych
dowodów dla kolejnych numerałów.

Cała procedura jest pierwotnie rekurencyjna,
ponieważ dla ustalonej $\delta$ stosuje jedynie
skończoną liczbę zagnieżdżonych pętli ograniczonych
i pierwotnie rekurencyjne operacje na kodach.
Uzyskany dowód jest początkowo dowodem w PA;
ponieważ $T\supseteq\mathrm{PA}$, można
jednolicie przeliczyć go na pełny kod dowodu $T$
z certyfikatami jego aksjomatów. PA formalizuje
tę konwersję. Otrzymujemy tezę.
\end{proof}

\begin{lemat}[Sformalizowana $\Sigma_1$-zupełność]
\label{lem:xi-formal-sigma1}
Dla ustalonej formuły $\Sigma_1$ postaci
$\sigma(\vec x)=\exists\vec u\,\delta(\vec x,\vec u)$,
gdzie $\delta$ jest $\Delta_0$, zachodzi
\[
\mathrm{PA}\vdash\forall\vec x\;
\bigl(\sigma(\vec x)\to
\operatorname{Bew}_T(\operatorname{inst}_\sigma(\vec x))\bigr).
\]
W szczególności dla zdania $\sigma$:
$\mathrm{PA}\vdash\sigma\to
\operatorname{Bew}_T(\overline{\ulcorner\sigma\urcorner})$.
\end{lemat}
\begin{proof}
Dla danych $\vec x,\vec u$ załóżmy
$\delta(\vec x,\vec u)$. Poprzedni lemat
dostarcza kod poprawnego dowodu
$\delta(\overline{\vec x},\overline{\vec u})$.
Kolejno dołączamy do niego wprowadzenia
kwantyfikatorów egzystencjalnych po wszystkich
zmiennych $\vec u$, otrzymując kod dowodu
$\sigma(\overline{\vec x})$. Ta operacja kodowa
jest pierwotnie rekurencyjna, a PA dowodzi jej
poprawności, zatem
\[
\mathrm{PA}\vdash
\forall\vec x\,\forall\vec u\;\left(
\delta(\vec x,\vec u)\to\exists p\,
\operatorname{Prf}_T(p,\operatorname{inst}_\sigma(\vec x))
\right).
\]
Wprowadzając kwantyfikatory egzystencjalne
po $\vec u$ po lewej stronie implikacji,
uzyskujemy $\sigma(\vec x)\to
\operatorname{Bew}_T(\operatorname{inst}_\sigma(\vec x))$.
To jest \emph{jedno} twierdzenie PA, również wtedy,
gdy wartości parametrów w modelu są niestandardowe.
\end{proof}

\begin{twierdzenie}[Warunki (D1)--(D3) dla $T$]
\label{thm:xi-hbl-verified}
Dla każdej efektywnie aksjomatyzowalnej teorii
$T\supseteq\mathrm{PA}$ ustalony standardowy
predykat dowodliwości z Lematu~\ref{res:7:7}
spełnia wszystkie warunki (D1)--(D3) z
Definicji~\ref{def:hilbert-bernays-loeb}.
\end{twierdzenie}
\begin{proof}
(D1) został udowodniony w Lemacie
\ref{lem:d1-provability}.
Warunek (D2) to Twierdzenie~\ref{thm:xi-d2}.
Dla (D3) ustalmy zdanie $\varphi$ i niech
$g:=\ulcorner\varphi\urcorner$. Ponieważ
$\operatorname{Bew}_T(x)$ jest $\Sigma_1$,
zdanie $\sigma:=\operatorname{Bew}_T(\overline g)$
ma postać $\Sigma_1$. Z Lematu
\ref{lem:xi-formal-sigma1} otrzymujemy
\[
\mathrm{PA}\vdash
\operatorname{Bew}_T(\overline g)\to
\operatorname{Bew}_T\left(
\overline{\ulcorner\operatorname{Bew}_T(\overline g)\urcorner}
\right).
\]
Pierwsza strona to $\Box\varphi$,
a druga to $\Box\Box\varphi$. Zatem (D3)
jest dowodliwe już w PA, więc także w $T$.
\end{proof}


\begin{przyklad}[Czym różni się dowód od jego kodu?]
Ustalmy dwie formalne dedukcje: pierwsza kończy się zdaniem $P$,
a druga kończy się zdaniem $P\to Q$. Ich teksty są obiektami
skończonymi, więc można je zakodować liczbami $q$ i $p$.
Metamatematycznie potrafimy połączyć teksty i dopisać regułę
modus ponens. To daje kod $m(p,q)$. Warunek (D2) mówi więcej:
\emph{ta sama operacja łączenia przechodzi dowód poprawności
wewnątrz arytmetyki}, ze zmiennymi $p,q$, a nie tylko po
wybraniu dwóch konkretnych liczb. Rozdzielenie tych poziomów
jest konieczne: skończona lista sprawdzonych przykładów nie
jest dowodem jednego zdania z dwoma kwantyfikatorami ogólnymi.
\end{przyklad}

\section{Czego nie wolno wyciągać z (D1)--(D3)?}
\begin{przyklad}[Trzy twierdzenia, trzy znaczenia]
Niech $\varphi$ oznacza konkretną równość
$\overline1+\overline1=\overline2$.
(D1) mówi, że mając formalny dowód tej równości,
można skonstruować dowód zdania $\Box\varphi$.
(D2) mówi, że teoria uzasadnia łączenie dowodów
implikacji i jej poprzednika. (D3) mówi,
że teoria potrafi formalnie potwierdzać istnienie
dowodu zdania o istnieniu dowodu.
Żaden z tych warunków nie jest zasadą
$\Box\varphi\to\varphi$ dla każdego $\varphi$.
\end{przyklad}

\begin{uwaga}[Dlaczego wybór predykatu jest istotny?]
Formuły zgodne na liczbach \emph{standardowych}
nie muszą być dowodliwie równoważne w PA.
Dlatego (D2) i (D3) są twierdzeniami
\emph{o ustalonej reprezentacji kodu dowodu},
a nie o dowolnej formule z odpowiednim zbiorem
wartości prawdziwych w $\mathfrak N$.
W szczególności predykat Rossera z
Definicji~\ref{def:rosser-predicate} nie otrzymuje
tych własności automatycznie.
\end{uwaga}